\documentclass[tikz,border=20pt]{standalone}
\usetikzlibrary{calc, arrows.meta}
\usepackage[UTF8]{ctex}

\begin{document}
\begin{tikzpicture}[
    iso/.style={very thick, draw=green!45!black},
    isof/.style={thick, draw=green!45!black},   % 淡态轮廓
    ghost/.style={semithick, draw=black!30, densely dashed},
    lab/.style={font=\sffamily\footnotesize, text=black!70, align=center},
    title/.style={font=\sffamily\bfseries\small, align=center},
    x=1cm, y=1cm,
]

% ============================================================
% 方案A：三态并排  before -> at -> after
% ============================================================
\begin{scope}[shift={(0,0)}]
  \node[title] at (2.4,2.1) {方案 A：三态并排（合并前 $\to$ 临界 $\to$ 合并后）};

  % before: 两个分离分量
  \begin{scope}[shift={(0,0)}, scale=0.62]
    \draw[iso] plot[smooth cycle, tension=0.75]
      coordinates {(-0.95,0.5)(-0.55,0.7)(-0.2,0.45)(-0.28,-0.05)
                   (-0.55,-0.5)(-1.0,-0.42)(-1.2,0.05)};
    \draw[iso] plot[smooth cycle, tension=0.75]
      coordinates {(0.95,0.5)(1.35,0.6)(1.6,0.1)(1.4,-0.42)
                   (0.95,-0.5)(0.55,-0.28)(0.6,0.15)};
  \end{scope}
  \node[lab] at (0,-1.55) {合并前\\$c=2$};

  \draw[-{Stealth[length=2mm]}, black!55, line width=0.8pt] (1.15,0) -- (1.75,0);

  % at: 相切临界
  \begin{scope}[shift={(2.9,0)}, scale=0.62]
    \draw[iso] plot[smooth cycle, tension=0.7]
      coordinates {(0,0.14)(-0.3,0.4)(-0.72,0.34)(-0.86,-0.05)
                   (-0.66,-0.4)(-0.26,-0.38)(0,-0.14)(-0.05,0)};
    \draw[iso] plot[smooth cycle, tension=0.7]
      coordinates {(0,0.14)(0.3,0.4)(0.72,0.34)(0.86,-0.05)
                   (0.66,-0.4)(0.26,-0.38)(0,-0.14)(0.05,0)};
    \node[circle, fill=red!80!black, draw=white, line width=0.4pt,
          inner sep=0pt, minimum size=3.2pt] at (0,0) {};
  \end{scope}
  \node[lab] at (2.9,-1.55) {临界（鞍点）\\相切一点};

  \draw[-{Stealth[length=2mm]}, black!55, line width=0.8pt] (4.05,0) -- (4.65,0);

  % after: 合并为单体
  \begin{scope}[shift={(5.8,0)}, scale=0.62]
    \draw[iso] plot[smooth cycle, tension=0.72]
      coordinates {(0,0.6)(0.75,0.5)(1.15,0.05)(0.85,-0.45)
                   (0,-0.6)(-0.85,-0.45)(-1.15,0.05)(-0.75,0.5)};
  \end{scope}
  \node[lab] at (5.8,-1.55) {合并后\\$c=1$};
\end{scope}

% ============================================================
% 方案B：临界态强调（相切点高亮 + 局部放大）
% ============================================================
\begin{scope}[shift={(0,-4.2)}]
  \node[title] at (2.4,2.1) {方案 B：临界态强调（相切点高亮 $+$ 局部放大）};

  \begin{scope}[shift={(1.6,0)}, scale=0.85]
    \draw[iso] plot[smooth cycle, tension=0.7]
      coordinates {(0,0.16)(-0.32,0.46)(-0.82,0.38)(-0.98,-0.05)
                   (-0.76,-0.46)(-0.3,-0.42)(0,-0.16)(-0.06,0)};
    \draw[iso] plot[smooth cycle, tension=0.7]
      coordinates {(0,0.16)(0.32,0.46)(0.82,0.38)(0.98,-0.05)
                   (0.76,-0.46)(0.3,-0.42)(0,-0.16)(0.06,0)};
    % 相切点高亮
    \draw[red!75!black, line width=1pt] (0,0) circle (0.16);
    \node[circle, fill=red!80!black, draw=white, line width=0.5pt,
          inner sep=0pt, minimum size=4pt] at (0,0) {};
  \end{scope}
  % 引出放大镜
  \draw[black!45, line width=0.7pt] (1.72,0.12) -- (3.7,1.0);
  \draw[black!45, line width=0.7pt] (1.72,-0.12) -- (3.7,-1.0);
  \draw[black!45, line width=0.9pt] (4.5,0) circle (1.05);
  \begin{scope}[shift={(4.5,0)}, scale=2.6]
    \draw[iso] plot[smooth, tension=0.9] coordinates {(-0.42,0.3)(-0.12,0.06)(0,0)};
    \draw[iso] plot[smooth, tension=0.9] coordinates {(-0.42,-0.3)(-0.12,-0.06)(0,0)};
    \draw[iso] plot[smooth, tension=0.9] coordinates {(0.42,0.3)(0.12,0.06)(0,0)};
    \draw[iso] plot[smooth, tension=0.9] coordinates {(0.42,-0.3)(0.12,-0.06)(0,0)};
    \node[circle, fill=red!80!black, draw=white, line width=0.3pt,
          inner sep=0pt, minimum size=2.4pt] at (0,0) {};
  \end{scope}
  \node[lab] at (4.5,-1.55) {相切点局部：\\两分量在此汇合};
\end{scope}

% ============================================================
% 方案C：上下夹心对比（穿越鞍点的连通性跳变）
% ============================================================
\begin{scope}[shift={(0,-8.4)}]
  \node[title] at (2.4,2.1) {方案 C：上下夹心（下=两分量,\ 中=临界,\ 上=已合并）};

  \begin{scope}[shift={(2.4,0)}, scale=0.9]
    % 下方 ghost：略低于鞍点，仍两分量
    \begin{scope}[shift={(0,-0.55)}, scale=0.92]
      \draw[ghost] plot[smooth cycle, tension=0.7]
        coordinates {(-0.18,0.12)(-0.4,0.34)(-0.78,0.28)(-0.9,-0.05)
                     (-0.7,-0.38)(-0.34,-0.34)(-0.18,-0.12)};
      \draw[ghost] plot[smooth cycle, tension=0.7]
        coordinates {(0.18,0.12)(0.4,0.34)(0.78,0.28)(0.9,-0.05)
                     (0.7,-0.38)(0.34,-0.34)(0.18,-0.12)};
    \end{scope}
    % 上方 ghost：略高于鞍点，合并单体
    \begin{scope}[shift={(0,0.6)}, scale=0.92]
      \draw[ghost] plot[smooth cycle, tension=0.72]
        coordinates {(0,0.42)(0.62,0.34)(0.82,0.0)(0.6,-0.34)
                     (0,-0.42)(-0.6,-0.34)(-0.82,0.0)(-0.62,0.34)};
    \end{scope}
    % 中间实线：临界相切态
    \draw[iso] plot[smooth cycle, tension=0.7]
      coordinates {(0,0.15)(-0.32,0.44)(-0.8,0.36)(-0.95,-0.05)
                   (-0.74,-0.44)(-0.3,-0.4)(0,-0.15)(-0.05,0)};
    \draw[iso] plot[smooth cycle, tension=0.7]
      coordinates {(0,0.15)(0.32,0.44)(0.8,0.36)(0.95,-0.05)
                   (0.74,-0.44)(0.3,-0.4)(0,-0.15)(0.05,0)};
    \node[circle, fill=red!80!black, draw=white, line width=0.5pt,
          inner sep=0pt, minimum size=3.6pt] at (0,0) {};
  \end{scope}
  \node[lab, text=black!45] at (4.0,0.6) {$f>f_s$：已合并 $c=1$};
  \node[lab, text=green!45!black] at (4.0,0.0) {$f=f_s$：临界相切};
  \node[lab, text=black!45] at (4.0,-0.55) {$f<f_s$：两分量 $c=2$};
\end{scope}

% ============================================================
% 第二组：本图真实情形 —— 鞍点上下都是“两个连通分量”（c:2->2 重连）
% ============================================================
\node[font=\sffamily\bfseries, text=blue!45!black, align=center]
  at (2.4,-11.4) {本图鞍点真实情形：上下都是 2 个分量（$c:2\!\to\!2$，分量接触后重新配对）};

% ---- 方案 D：接触—重连三态（下2分量 -> 相切重连 -> 上2分量，配对改变）----
\begin{scope}[shift={(0,-14.2)}]
  \node[title] at (2.4,1.9) {方案 D：接触—重连三态（$c$ 不变，配对重组）};

  % 下方：两分量“横向相邻”
  \begin{scope}[shift={(0,0)}, scale=0.6]
    \draw[iso] plot[smooth cycle, tension=0.75]
      coordinates {(-1.15,0.45)(-0.7,0.62)(-0.42,0.2)(-0.55,-0.3)
                   (-1.0,-0.5)(-1.45,-0.28)(-1.5,0.15)};
    \draw[iso] plot[smooth cycle, tension=0.75]
      coordinates {(1.15,0.45)(1.6,0.6)(1.6,0.05)(1.35,-0.42)
                   (0.9,-0.5)(0.45,-0.25)(0.55,0.22)};
  \end{scope}
  \node[lab] at (0,-1.5) {下方 $c=2$\\左右两分量};

  \draw[-{Stealth[length=2mm]}, black!55, line width=0.8pt] (1.35,0) -- (1.95,0);

  % 临界：两分量在鞍点相切于一点
  \begin{scope}[shift={(3.1,0)}, scale=0.6]
    \draw[iso] plot[smooth cycle, tension=0.7]
      coordinates {(0,0.12)(-0.34,0.42)(-0.82,0.34)(-0.96,-0.08)
                   (-0.7,-0.44)(-0.28,-0.4)(0,-0.12)(-0.05,0)};
    \draw[iso] plot[smooth cycle, tension=0.7]
      coordinates {(0,0.12)(0.34,0.42)(0.82,0.34)(0.96,-0.08)
                   (0.7,-0.44)(0.28,-0.4)(0,-0.12)(0.05,0)};
    \node[circle, fill=red!80!black, draw=white, line width=0.4pt,
          inner sep=0pt, minimum size=3.2pt] at (0,0) {};
  \end{scope}
  \node[lab] at (3.1,-1.5) {鞍点：相切于一点\\（重连临界）};

  \draw[-{Stealth[length=2mm]}, black!55, line width=0.8pt] (4.45,0) -- (5.05,0);

  % 上方：仍两分量，但改为“上下配对”（重连后连接方式变）
  \begin{scope}[shift={(6.3,0)}, scale=0.6]
    \draw[iso] plot[smooth cycle, tension=0.75]
      coordinates {(0,1.15)(0.5,0.95)(0.62,0.55)(0.3,0.42)
                   (-0.3,0.5)(-0.55,0.9)(-0.35,1.15)};
    \draw[iso] plot[smooth cycle, tension=0.75]
      coordinates {(0,-1.15)(0.55,-0.95)(0.6,-0.5)(0.28,-0.42)
                   (-0.35,-0.5)(-0.6,-0.9)(-0.3,-1.15)};
  \end{scope}
  \node[lab] at (6.3,-1.5) {上方 $c=2$\\两分量重组};
\end{scope}

% ---- 方案 E：重连局部放大（聚焦“X 形四臂”交换）----
\begin{scope}[shift={(0,-18.6)}]
  \node[title] at (2.4,1.9) {方案 E：重连局部放大（四臂在鞍点交换连接）};

  \begin{scope}[shift={(1.5,0)}, scale=0.8]
    \draw[iso] plot[smooth cycle, tension=0.7]
      coordinates {(0,0.14)(-0.34,0.44)(-0.84,0.36)(-0.98,-0.06)
                   (-0.72,-0.44)(-0.3,-0.4)(0,-0.14)(-0.05,0)};
    \draw[iso] plot[smooth cycle, tension=0.7]
      coordinates {(0,0.14)(0.34,0.44)(0.84,0.36)(0.98,-0.06)
                   (0.72,-0.44)(0.3,-0.4)(0,-0.14)(0.05,0)};
    \draw[red!75!black, line width=1pt] (0,0) circle (0.15);
    \node[circle, fill=red!80!black, draw=white, line width=0.5pt,
          inner sep=0pt, minimum size=4pt] at (0,0) {};
  \end{scope}
  \draw[black!45, line width=0.7pt] (1.6,0.12) -- (3.7,1.0);
  \draw[black!45, line width=0.7pt] (1.6,-0.12) -- (3.7,-1.0);
  \draw[black!45, line width=0.9pt] (4.5,0) circle (1.05);
  % 放大：四臂 X 形，配彩色区分“重连前后归属”
  \begin{scope}[shift={(4.5,0)}, scale=2.6]
    \draw[iso] plot[smooth, tension=0.9] coordinates {(-0.42,0.3)(-0.12,0.06)(0,0)};
    \draw[iso] plot[smooth, tension=0.9] coordinates {(0.42,0.3)(0.12,0.06)(0,0)};
    \draw[iso] plot[smooth, tension=0.9] coordinates {(-0.42,-0.3)(-0.12,-0.06)(0,0)};
    \draw[iso] plot[smooth, tension=0.9] coordinates {(0.42,-0.3)(0.12,-0.06)(0,0)};
    % 重连方向示意：上下配 / 左右配
    \draw[blue!55!black, densely dashed, line width=0.5pt] (-0.28,0.2) .. controls (0,0.05) .. (0.28,0.2);
    \draw[orange!75!black, densely dashed, line width=0.5pt] (-0.28,-0.2) .. controls (0,-0.05) .. (0.28,-0.2);
    \node[circle, fill=red!80!black, draw=white, line width=0.3pt,
          inner sep=0pt, minimum size=2.4pt] at (0,0) {};
  \end{scope}
  \node[lab] at (4.5,-1.5) {四臂在鞍点汇于一点\\穿越后连接方式交换};
\end{scope}

% ---- 方案 F：上下夹心（下2分量 / 中临界 / 上2分量，形态重组）----
\begin{scope}[shift={(0,-23.0)}]
  \node[title] at (2.4,1.9) {方案 F：上下夹心（$c$ 恒为 2，仅分量形态重组）};

  \begin{scope}[shift={(2.4,0)}, scale=0.9]
    % 下方 ghost：左右两分量
    \begin{scope}[shift={(0,-0.62)}, scale=0.9]
      \draw[ghost] plot[smooth cycle, tension=0.72]
        coordinates {(-0.2,0.1)(-0.42,0.32)(-0.8,0.26)(-0.92,-0.06)
                     (-0.72,-0.36)(-0.36,-0.32)(-0.2,-0.1)};
      \draw[ghost] plot[smooth cycle, tension=0.72]
        coordinates {(0.2,0.1)(0.42,0.32)(0.8,0.26)(0.92,-0.06)
                     (0.72,-0.36)(0.36,-0.32)(0.2,-0.1)};
    \end{scope}
    % 上方 ghost：上下两分量（重连后）
    \begin{scope}[shift={(0,0.66)}, scale=0.82]
      \draw[ghost] plot[smooth cycle, tension=0.72]
        coordinates {(0,0.5)(0.4,0.34)(0.42,0.05)(0.16,-0.04)(-0.16,-0.02)(-0.42,0.06)(-0.4,0.34)};
    \end{scope}
    \begin{scope}[shift={(0,-0.02)}]
    \end{scope}
    % 中间实线：临界相切态
    \draw[iso] plot[smooth cycle, tension=0.7]
      coordinates {(0,0.14)(-0.34,0.42)(-0.82,0.34)(-0.96,-0.06)
                   (-0.72,-0.42)(-0.3,-0.38)(0,-0.14)(-0.05,0)};
    \draw[iso] plot[smooth cycle, tension=0.7]
      coordinates {(0,0.14)(0.34,0.42)(0.82,0.34)(0.96,-0.06)
                   (0.72,-0.42)(0.3,-0.38)(0,-0.14)(0.05,0)};
    \node[circle, fill=red!80!black, draw=white, line width=0.5pt,
          inner sep=0pt, minimum size=3.6pt] at (0,0) {};
  \end{scope}
  \node[lab, text=black!45] at (4.1,0.66) {$f>f_s$：仍 $c=2$（重组）};
  \node[lab, text=green!45!black] at (4.1,0.0) {$f=f_s$：相切重连};
  \node[lab, text=black!45] at (4.1,-0.62) {$f<f_s$：$c=2$（左右）};
\end{scope}

\end{tikzpicture}
\end{document}